\subsection{GROUP GERMS}

DEFINITION 5. A Lie group germ over K is a system (G, e, θ, m) satisfying the following conditions:

\begin{enumerate}
\def\labelenumi{(\roman{enumi})}
\item
  G is an analytic manifold over K;
\item
  \(e \in G;\)
\item
  \(\theta\) is an analytic mapping of G into G;
\item
  \(m\) is an analytic mapping of an open subset \(\Omega\) of \(\mathbf{G} \times \mathbf{G}\) into \(\mathbf{G}\) ;
\item
  for all \(g \in G\) , \((e, g) \in \Omega\) , \((g, e) \in \Omega\) , \(m(e, g) = m(g, e) = g\) ;
\end{enumerate}

\[
g \in G, (g, \theta (g)) \in \Omega , (\theta (g), g) \in \Omega , m (g, \theta (g)) = m (\theta (g), g) = e;
\]

\begin{enumerate}
\def\labelenumi{(\roman{enumi})}
\setcounter{enumi}{6}
\tightlist
\item
  if \(g, h, k\) are elements of \(G\) such that \((g, h) \in \Omega\) , \((h, k) \in \Omega\) , \((m(g, h), k) \in \Omega\) , \((g, m(h, k)) \in \Omega\) , then \(m(m(g, h), k) = m(g, m(h, k))\) .
\end{enumerate}

e is called the identity element of the group germ. We often write gh instead of \(m(g, h)\) and (by an abuse of notation) \(g^{-1}\) instead of \(\theta(g)\) .

A Lie group G is a Lie group germ with the obvious choice of \(e, \theta, m\) . Let G be a Lie group germ. Then \(ee^{-1} = e\) , that is

\[
e ^ {- 1} = e.\tag{8}
\]

For all \(g\in G\)

\[
g = e g = ((g ^ {- 1}) ^ {- 1} g ^ {- 1}) g = (g ^ {- 1}) ^ {- 1} (g ^ {- 1} g) = (g ^ {- 1}) ^ {- 1} e,
\]

that is

\[
(g ^ {- 1}) ^ {- 1} = g.\tag{9}
\]

A subset of G invariant under the mapping \(g \mapsto g^{-1}\) is called symmetric. The manifold G, with the point e, the mapping \(g \mapsto g^{-1}\) and the mapping \((g, h) \mapsto hg\) is a Lie Group germ G \(^{\nu}\) called the opposite of G.

The Lie group germ G is called commutative if, for all \((g,h)\in \mathbf{G}\times \mathbf{G}\) such that \(gh\) is defined, \(hg\) is defined and equal to \(gh\) .

Let \(G\) be a Lie group germ. The set of \((g, h) \in G \times G\) such that \(gh\) is defined is a neighbourhood of \((e, e)\) . On the other hand, the mappings \((g, h) \mapsto gh\) and \(g \mapsto g^{-1}\) are continuous. Hence \((gh)k = g(hk)\) for \(g, h, k\) sufficiently close to \(e\) . Similarly, \((h^{-1}g^{-1})(gh) = h^{-1}(eh) = h^{-1}h = e\) for \(g, h\) sufficiently close to \(e\) , whence multiplying on the right by \((gh)^{-1}\) ,

\begin{enumerate}
\def\labelenumi{(\arabic{enumi})}
\setcounter{enumi}{9}
\tightlist
\item
\end{enumerate}

\((gh)^{-1} = h^{-1}g^{-1}\) for \(g,h\) sufficiently close to \(e\) .

PROPOSITION 20. Let G be a Lie group germ and \(g \in G\) . There exist an open neighbourhood U of e and an open neighbourhood V of g with the following properties:

\begin{enumerate}
\def\labelenumi{(\alph{enumi})}
\item
  ug is defined for all \(u \in U\) ;
\item
  \(vg^{-1}\) is defined for all \(v \in V\) ;
\item
  the mappings \(u \mapsto ug\) , \(v \mapsto vg^{-1}\) are inverse analytic isomorphisms of one another of U onto V and of V onto U.
\end{enumerate}

As the set of definition of the product is open in \(\mathbf{G} \times \mathbf{G}\) , there exist an open neighbourhood U of \(e\) and an open neighbourhood V of \(g\) with properties (a) and (b). Let \(\eta(u) = ug\) for \(u \in \mathrm{U}\) , \(\eta'(v) = vg^{-1}\) for \(v \in \mathrm{V}\) . By shrinking U and V, it can be assumed that \((ug)g^{-1} = u\) and \((vg^{-1})g = v\) for \(u \in \mathrm{U}\) and \(v \in \mathrm{V}\) . Then \(\eta\) and \(\eta'\) are injections. By shrinking U further, it can be assumed that \(\eta(\mathrm{U}) \subset \mathrm{V}\) . Then \(\eta'(\mathrm{V}) \supset \mathrm{U}\) and \(\eta(\mathrm{U})\) is the inverse image of U under \(\eta'\) and hence is an open neighbourhood of \(g\) in V. Replacing V by \(\eta(\mathrm{U})\) , we finally arrive at the situation where \(\eta\) and \(\eta'\) are analytic inverse bijections.

Let \(G_{1}, G_{2}\) be two Lie group germs with identity elements \(e_{1}, e_{2}\) . A mapping f of \(G_{1}\) into \(G_{2}\) is called a morphism if f satisfies the following conditions:

\begin{enumerate}
\def\labelenumi{(\roman{enumi})}
\item
  \(f\) is analytic;
\item
  \(f(e_{1}) = e_{2};\)
\item
  if \(g, h\) are elements of \(\mathbf{G}_1\) such that \(gh\) is defined, then \(f(g)f(h)\) is defined and is equal to \(f(gh)\) .
\end{enumerate}

Let \(g \in G_1\) . As \(gg^{-1}\) is defined and is equal to \(e_1, f(g)f(g^{-1})\) is defined and is equal to \(e_2\) and hence

\[
f (g) ^ {- 1} = f (g) ^ {- 1} \left(f (g) f \left(g ^ {- 1}\right)\right) = \left(f (g) ^ {- 1} f (g)\right) f \left(g ^ {- 1}\right)
\]

that is

\[
f (g) ^ {- 1} = f (g ^ {- 1}).\tag{11}
\]

The composition of two morphisms is a morphism.

If \(f: \mathbf{G}_1 \to \mathbf{G}_2\) and \(f': \mathbf{G}_2 \to \mathbf{G}_1\) are two inverse morphisms of one another, they are isomorphisms (using in particular formula (11)).

Let \(G_{1}\) , \(G_{2}\) be two Lie group germs, and f a mapping of \(G_{1}\) into \(G_{2}\) satisfying conditions (ii) and (iii) above, which is analytic in an open neighbourhood of \(e_{1}\) . Using Proposition 20, it can be proved as in Proposition 4 that f is a morphism.

Let \((\mathbf{G}, e, \theta, m)\) be a Lie group germ and \(\Omega\) the set of definition of \(m\) . Let \(\mathbf{H}\) be a submanifold of \(\mathbf{G}\) containing \(e\) , which is stable under \(\theta\) . Suppose that the set \(\Omega_1\) of \((x, y) \in \Omega \cap (\mathbf{H} \times \mathbf{H})\) such that \(m(x, y) \in \mathbf{H}\) is open in \(\mathbf{H} \times \mathbf{H}\) . Then \((\mathbf{H}, e, \theta | \mathbf{H}, m | \Omega_1)\) is a Lie group germ. Such a Lie group germ is called a Lie subgroup germ of \(\mathbf{G}\) . The canonical injection of \(\mathbf{H}\) into \(\mathbf{G}\) is a morphism. If \(f: \mathbf{L} \to \mathbf{G}\) is a morphism of Lie group germs such that \(f(\mathbf{L}) \subset \mathbf{H}\) , then \(f: \mathbf{L} \to \mathbf{H}\) is a morphism of Lie group germs.

Suppose that K is of characteristic 0. If we replace the hypothesis that H is a submanifold of G by the hypothesis that H is a quasi-submanifold of G, the results of the above paragraph remain true (cf.~Differentiable and Analytic Manifolds, R, 5.8.5). H is then called a Lie quasi-subgroup germ of G.

If G is a Lie group germ with identity element e, every symmetric open neighbourhood of e in G is a Lie subgroup germ of G. (This applies in particular when G is a Lie group.) Let H be a Lie subgroup germ of G; if H is a neighbourhood of e in G, then H is open in G by Proposition 20.

The product Lie group germ of a finite number of Lie group germs is defined in an obvious way.

PROPOSITION 21. Let G, H be two Lie group germs and \(\phi\) a morphism of G into H. The following conditions are equivalent:

\begin{enumerate}
\def\labelenumi{(\roman{enumi})}
\item
  \(\phi\) is étale at \(e\) ;
\item
  there exist open Lie subgroup germs G', H' of G, H such that \(\phi|G'\) is an isomorphism of G' onto H'.
\end{enumerate}

The implication \(\langle \mathrm{ii}\rangle \Rightarrow \langle \mathrm{i}\rangle\) is obvious. Suppose that \(\phi\) is étale at \(e\) . There exists an open Lie subgroup germ \(\mathbf{G}_1\) of \(\mathbf{G}\) such that \(\phi (\mathbf{G}_1)\) is open in \(\mathbf{H}\) and \(\phi |\mathbf{G}_1\) is an isomorphism of the manifold \(\mathbf{G}_1\) onto the manifold \(\phi (\mathbf{G}_1)\) . Then there exists an open Lie subgroup germ \(\mathbf{G}'\) of \(\mathbf{G}_1\) such that the product in \(\mathbf{G}\) of two elements of \(\mathbf{G}'\) is always defined and belongs to \(\mathbf{G}_1\) . If \(g, g'\) are elements of \(\mathbf{G}'\) such that \(\mathrm{gg}^{\prime}\in \mathrm{G}'\) , then \(\phi (g)\phi (g') = \phi (gg')\in \phi (\mathbf{G}')\) ; if \(g, g'\) are elements of \(\mathbf{G}'\) such that \(gg^{\prime}\in G_1 - G'\) , then

\[
\phi (g) \phi (g ^ {\prime}) = \phi (g g ^ {\prime}) \in \phi (\mathrm{G} _ {1}) - \phi (\mathrm{G} ^ {\prime}).
\]

Hence \(\phi |G^{\prime}\) is an isomorphism of the Lie group germ \(G^{\prime}\) onto the open Lie subgroup germ \(\phi (G^{\prime})\) of H.

If the conditions of Proposition 21 hold, G and H are called locally isomorphic.

PROPOSITION 22. Let H be a Lie group, U a Lie subgroup germ of H and N the set of \(g \in \mathrm{H}\) such that U and \(g\mathrm{Ug}^{-1}\) have the same germ at \(e\) (General Topology, Chapter I, § 6, no. 10). Then \(N\) is a subgroup of H containing U. There exists one and only one analytic manifold structure on \(\mathbb{N}\) with the following properties:

\begin{enumerate}
\def\labelenumi{(\roman{enumi})}
\item
  N with this structure is a Lie group;
\item
  U is an open submanifold of N;
\item
  the canonical injection of \(\mathbf{N}\) into \(\mathbf{H}\) is an immersion.
\end{enumerate}

Clearly \(\mathbf{N}\) is a subgroup of \(\mathbf{H}\) . If \(g \in \mathbf{U}\) , then \(ge \in \mathbf{U}\) and \(geg^{-1} \in \mathbf{U}\) , hence \(gu \in \mathbf{U}\) and \(gug^{-1} \in \mathbf{U}\) for \(u\) sufficiently close to \(e\) in \(\mathbf{U}\) and hence the germ of \(g\mathrm{U}g^{-1}\) at \(e\) is contained in that of \(\mathbf{U}\) ; exchanging \(g\) and \(g^{-1}\) , we see that the germs of \(g\mathrm{U}g^{-1}\) and \(\mathbf{U}\) at \(e\) are equal. Hence \(\mathbf{U} \subset \mathbf{N}\) .

Let \(\mathbf{V}\) be an open neighbourhood of \(e\) in \(\mathbf{U}\) such that \(\mathrm{V} = \mathrm{V}^{-1}\) , \(\mathrm{V}^2 \subset \mathrm{U}\) . Conditions (i), (ii), (iii) of Proposition 18 of no. 9 (where \(\mathrm{G}\) is replaced by \(\mathrm{N}\) ) are satisfied. Hence there exists an analytic manifold structure on \(\mathbf{N}\) with the following properties: \((\alpha)\) \(\mathbf{N}\) with this structure is a Lie group; \((\beta)\) \(\mathbf{V}\) is open in \(\mathbf{N}\) ; \((\gamma)\) the manifold structures on \(\mathbf{N}\) and \(\mathbf{U}\) induce the same structure on \(\mathbf{V}\) . Since \(\mathbf{V}\) is a submanifold of \(\mathbf{H}\) , the canonical injection of \(\mathbf{N}\) into \(\mathbf{H}\) is an immersion at \(e\) and hence at every point of \(\mathbf{N}\) . Let \(u \in \mathbf{U}\) . There exists an open neighbourhood \(V'\) of e in V such that the mapping \(v \mapsto vu\) is an analytic isomorphism of \(V'\) onto an open neighbourhood of u in U (Proposition 20) and at the same time onto an open neighbourhood of u in N. Hence U is open in N and the identity mapping of U is an isomorphism with the given manifold structure on U and the open submanifold structure on N; in other words, U is an open submanifold of N.

Finally, we consider an analytic manifold structure on N with properties (i) and (ii) of the proposition and let \(N^{*}\) be the Lie group thus obtained. Then the identity mapping of N into \(N^{*}\) is étale at e and hence a Lie group isomorphism. This proves the uniqueness assertion of the proposition.

Let H be a Lie group, U a Lie quasi-subgroup germ of H and N the set of \(g \in H\) such that U and \(gUg^{-1}\) have the same germ at e. If K is of characteristic 0, there exists on G one and only one manifold structure with properties (i) and (ii) of Proposition 22. The proof is the same as for Proposition 22.

COROLLARY. Preserving the notation of Proposition 22, let G be the subgroup of H generated by U. Then G is an open subgroup of N. There exists one and only one Lie group structure on G such that U is an open submanifold of G and the canonical injection of G into H is an immersion.

Remark. Preserving the notation of Proposition 22 and its corollary, suppose that K is of characteristic 0, that H is finite-dimensional and that the topology on U admits a countable base. Even with all these hypotheses it is possible that G is not closed in H (Exercise 3). But, if G is closed, G is a Lie subgroup of H. For the mapping \((g,h)\mapsto gh\) is a law of analytic left operation of G on H. The orbit of e is G. Our assertion then follows from Propositions 2 (iv) and 14 (iii).

